\paragraph{Normalization of alpha and its signed configuration.}
\label{inc-ampl-clausura}
The coordinate of $\alpha$ is determined along the positional route by
composing regional blocks and the dodecaphase record. Before
positional division, the oriented combination has components
$Z_j=P_j+E_j-\Phi_j-K_j^{\mathrm{per}}$, and the balance with incoming quotient
is $s_j=Z_j+c_{j+1}$. Normalization satisfies
\begin{equation}
 A_j=s_j-1000c_j=Z_j+c_{j+1}-1000c_j,\qquad 0\leq A_j<1000.
 \label{inc-ampl-normalizacion}
\end{equation}
This identity simultaneously preserves the normalized block and the
contribution of the transported quotient. The limit of compatible
prefixes determines the scalar closure coordinate. The
signed configuration, for its part, preserves which positions in the
combination before normalization exhibit negative defect.
In the initial window, that configuration is the vector $Z_0$ of
\eqref{exc:precarry}, and its negative support is $H^-_\alpha$.

Two contents of the same operation are thus defined. Positional
normalization determines a real coordinate through its compatible
remainders and quotients; negative support determines a
configuration of positions in the twelve-component cycle. The sign
refers to the balance of records before normalization. Its incidence
reading remains available because the state preserves that balance
alongside the normalized result. This articulation corresponds to the
proved positional route; comparison with the analytic vacancy
operators preserves the precise scope established in the chapter on
$\alpha$.
